%% =====================================================================
%%  9 Conclusion
%% =====================================================================
%%  GUIDE
%%   - One numbered sentence per question, with the same numbers as the
%%     summary and the results section (identical digits -- judges notice).
%%   - One closing sentence: what the result means beyond this problem.
%%   - No new material, no new citations, no speculation.
%% =====================================================================
\section{Conclusion}
\label{sec:conclusion}

\begin{enumerate}[label=(\arabic*),leftmargin=*]
  \item A \SI{10}{cm} breach in a \SI{628}{m^3} module at \SI{101}{kPa}
        drops the internal pressure to $0.3\,p_0$ in \SI{286}{s}, and the
        time scales as $t_{0.3}\propto d^{-1.97}$ between \SI{5}{cm} and
        \SI{20}{cm}.
  \item A crew needing \SI{120}{s} to seal the breach keeps a margin of
        \SI{166}{s} at \SI{10}{cm}; the margin vanishes at
        \SI{15.4}{cm}, above which the module must be abandoned.
  \item The answer is uncertain at the $\pm\SI{12}{\percent}$ level, with
        \SI{10}{\percent} of that coming from the hole diameter and
        \SI{6}{\percent} from the discharge coefficient; measuring the
        hole is worth more than refining the model.
\end{enumerate}

Because the two governing groups are $\gamma$ and $\Pi_T=\tau/\tau_w$, the
same numbers transfer to any rigid volume venting into vacuum once its
size and hole diameter are expressed through $\tau$ --- the result is a
design rule for breach response, not a single-case calculation.
